\documentclass[12pt]{article}
\usepackage{geometry}
\geometry{margin=1in}

% ----------------------------------------------------------
%  Linux Penguin Theme — Based on Recommended Colors
% ----------------------------------------------------------

% --- COLOR SETUP ---
\usepackage{xcolor}

\definecolor{cream}{RGB}{246,242,219}
\definecolor{teamred}{RGB}{209,57,64}
\definecolor{navy}{RGB}{40,46,87}
\definecolor{softwhite}{RGB}{250,250,230}
\definecolor{accentyellow}{RGB}{240,200,60}

% Extra utility colors
\definecolor{muted}{RGB}{120,125,140}
\definecolor{darkcream}{RGB}{230,225,205}

% ----------------------------------------------------------
%  FONT SETUP (Optional but Recommended)
%  Works best with XeLaTeX or LuaLaTeX
% ----------------------------------------------------------
\usepackage{fontspec}

% Main document font (Google fonts — upload to Overleaf)
% If you don't upload fonts, comment these two lines.
\setmainfont{SourceSans3}[
  Path=./,
  ItalicFont  = *-Italic,
  BoldFont    = *-Bold
]

\newfontfamily\headingfont{ScienceGothic}[
  Path=./,
]

% Font Awesome for minimal icons
\usepackage{fontawesome5}

% ----------------------------------------------------------
%  SECTION STYLE
% ----------------------------------------------------------
\usepackage{titlesec}

% Section
\titleformat{\section}
  {\Large\headingfont\bfseries\color{navy}}
  {\thesection}
  {1em}
  {\rule{\linewidth}{1pt}\vspace{0.5em}\\}

% Subsection
\titleformat{\subsection}
  {\large\headingfont\bfseries\color{teamred}}
  {\thesubsection}
  {0.8em}
  {}

% Subsubsection
\titleformat{\subsubsection}
  {\bfseries\color{navy}}
  {\thesubsubsection}
  {0.6em}
  {}

% ----------------------------------------------------------
%  PAGE STYLE
% ----------------------------------------------------------
\usepackage{fancyhdr}
\pagestyle{fancy}

\fancyhf{} % clear all

\lhead{\textcolor{navy}{\headingfont \courseTitle}}
\rhead{\textcolor{teamred}{\thepage}}
\renewcommand{\headrule}{\color{teamred}\hrule height 1pt}

% ----------------------------------------------------------
%  HYPERLINK STYLE
% ----------------------------------------------------------
\usepackage{hyperref}

\hypersetup{
  colorlinks = true,
  linkcolor  = teamred,
  urlcolor   = navy,
  citecolor  = accentyellow
}

% ----------------------------------------------------------
%  CUSTOM BOX ENVIRONMENTS
% ----------------------------------------------------------
\usepackage{tcolorbox}
\tcbuselibrary{skins,breakable}

% Info box
\newtcolorbox{infobox}{
  enhanced,
  breakable,
  colback   = softwhite,
  colframe  = navy,
  coltext   = navy,
  boxrule   = 1pt,
  arc       = 4pt,
  left=6pt, right=6pt, top=6pt, bottom=6pt
}

% Warning box
\newtcolorbox{warningbox}{
  enhanced,
  breakable,
  colback   = accentyellow!30,
  colframe  = accentyellow!80!navy,
  coltext   = navy,
  boxrule   = 1pt,
  arc       = 4pt,
  left=6pt, right=6pt, top=6pt, bottom=6pt
}

% Highlight box
\newtcolorbox{highlightbox}{
  enhanced,
  breakable,
  colback   = teamred!10,
  colframe  = teamred,
  coltext   = navy,
  boxrule   = 1pt,
  arc       = 4pt,
  left=6pt, right=6pt, top=6pt, bottom=6pt
}

% ----------------------------------------------------------
%  ITEMIZE / ENUMERATE STYLE
% ----------------------------------------------------------
\usepackage{enumitem}

\setlist[itemize]{
  label=\textcolor{teamred}{\large\textbullet},
  leftmargin=1.2em,
  itemsep=4pt
}

\setlist[enumerate]{
  label=\textcolor{navy}{\arabic*.},
  leftmargin=1.4em,
  itemsep=4pt
}

% ----------------------------------------------------------
%  TITLE STYLE
% ----------------------------------------------------------
\usepackage{titling}

\pretitle{
  \begin{center}
  \headingfont\bfseries\LARGE\color{navy}
}
\posttitle{
  \end{center}
}

\preauthor{
  \begin{center}
  \color{teamred}
}
\postauthor{
  \end{center}
}

% ----------------------------------------------------------
%  OPTIONAL — SOFT PAGE BACKGROUND
%  Comment out if you want white pages
% ----------------------------------------------------------
\usepackage{pagecolor}
\pagecolor{cream}
\color{navy}

% ----------------------------------------------------------
% CUSTOM COMMANDS FOR TEMPLATE
% ----------------------------------------------------------
\newcommand{\courseTitle}{Scheduling with systemd on Linux}
\newcommand{\courseDate}{28-02-2026}
\newcommand{\authorName}{Khaled Mahfouz}
\newcommand{\contactInfo}{\texttt{@khaledmhfz}}

% Minimal icon commands
\newcommand{\penguinicon}{\faLinux}
\newcommand{\systemdicon}{\faCogs}
\newcommand{\serviceicon}{\faPlay}
\newcommand{\timericon}{\faClock}
\newcommand{\uniticon}{\faFile}
\newcommand{\controlicon}{\faSlidersH}
\newcommand{\warningicon}{\faExclamationTriangle}
\newcommand{\tipsicon}{\faLightbulb}
\newcommand{\arrowicon}{\faArrowRight}
\newcommand{\bookicon}{\faBook}
\newcommand{\rocketicon}{\faRocket}

% ----------------------------------------------------------
% END OF THEME
% ----------------------------------------------------------

\begin{document}

% Title with minimal space
\begin{center}
\headingfont\Huge\bfseries\color{navy}
\systemdicon\ \courseTitle\\[0.5em]
\Large\color{teamred}
\courseDate
\end{center}

\begin{infobox}
Alright team \faUserTie\ — this is the session where you level up from "I run Linux" to:

\begin{center}
\large\textbf{"I control what runs on Linux… and when it runs."}
\end{center}
\end{infobox}

We're diving into:

\begin{itemize}
    \item \texttt{systemd} services
    \item Basic unit files
    \item \texttt{systemd} timers (modern scheduling)
\end{itemize}

By the end, you'll:

\begin{itemize}
    \item Understand how your system boots and manages services
    \item Create your own service units
    \item Schedule tasks like a pro (without cron)
    \item Override and customize services safely
\end{itemize}

\vspace{1em}
\begin{center}
\color{teamred}\rule{0.8\textwidth}{1pt}
\end{center}

\section*{1 \faInfoCircle\ What is systemd? (And Why You Should Care)}

Before scheduling anything, you need to understand what's running your system.

On modern Debian-based systems:

\begin{itemize}
    \item \textbf{Debian}
    \item \textbf{Ubuntu}
    \item \textbf{Linux Mint}
\end{itemize}

The init system is:

\subsection*{\faBrain\ \texttt{systemd}}

An \textit{init system} is the first userspace process started by the kernel.

Check it:

\begin{highlightbox}
\texttt{\$ ps -p 1 -o comm=}
\end{highlightbox}

You'll see:

\begin{highlightbox}
\texttt{systemd}
\end{highlightbox}

\subsection*{\faQuestionCircle\ What Does systemd Do?}

\begin{itemize}
    \item Starts services during boot
    \item Manages background processes (daemons)
    \item Handles logging (via \texttt{journald})
    \item Mounts filesystems
    \item Manages network services
    \item Handles timers (our main topic today)
\end{itemize}

\begin{infobox}
If Linux is a city \faCity\\
\texttt{systemd} is the mayor.
\end{infobox}

\subsection*{\bookicon\ Helpful Reading}

\begin{itemize}
    \item \href{https://www.geeksforgeeks.org/linux-unix/linux-systemd-and-its-components/}{\texttt{geeksforgeeks-linux-systemd-and-its-components}}
    \item \href{https://www.geeksforgeeks.org/linux-unix/an-introduction-to-systemd-and-its-role-in-the-boot-process/}{\texttt{geeksforgeeks-introduction-to-systemd}}
\end{itemize}

\section*{2 \uniticon\ Understanding Units in systemd}

Everything in systemd is a \textbf{Unit}.

A unit is a configuration file that tells systemd what to manage.

\subsection*{Common Unit Types}

\begin{infobox}
\begin{tabular}{lll}
\textbf{Type} & \textbf{Extension} & \textbf{Purpose} \\
Service & \texttt{.service} & Background services \\
Timer & \texttt{.timer} & Scheduling \\
Mount & \texttt{.mount} & Mount points \\
Target & \texttt{.target} & Group of units \\
Socket & \texttt{.socket} & Socket activation \\
\end{tabular}
\end{infobox}

Example:

\begin{highlightbox}
\texttt{apache2.service}\\
\texttt{ssh.service}\\
\texttt{cron.service}
\end{highlightbox}

To list all units:

\begin{highlightbox}
\texttt{\$ systemctl list-units}
\end{highlightbox}

To list all installed units:

\begin{highlightbox}
\texttt{\$ systemctl list-unit-files}
\end{highlightbox}

\section*{3 \faDownload\ Installing Apache (Service Example)}

Let's install something real:

\begin{highlightbox}
\texttt{\$ sudo apt update}\\
\texttt{\$ sudo apt install apache2}
\end{highlightbox}

This installs:

\begin{highlightbox}
\texttt{apache2.service}
\end{highlightbox}

Now check its status:

\begin{highlightbox}
\texttt{\$ systemctl status apache2}
\end{highlightbox}

You'll see:

\begin{itemize}
    \item Loaded
    \item Active
    \item Main PID
    \item Logs
\end{itemize}

\section*{4 \controlicon\ Managing Services with systemctl}

This is daily Linux admin stuff.

\subsection*{\faPlay\ Start a Service}

\begin{highlightbox}
\texttt{\$ sudo systemctl start apache2}
\end{highlightbox}

\subsection*{\faStop\ Stop a Service}

\begin{highlightbox}
\texttt{\$ sudo systemctl stop apache2}
\end{highlightbox}

\subsection*{\faRedo\ Restart}

\begin{highlightbox}
\texttt{\$ sudo systemctl restart apache2}
\end{highlightbox}

\subsection*{\faSync\ Reload vs Restart (Important!)}

\begin{itemize}
    \item \textbf{\texttt{reload}}
    \begin{itemize}
        \item Reloads configuration
        \item Does NOT stop process
        \item Keeps connections alive
    \end{itemize}
    
    \item \textbf{\texttt{restart}}
    \begin{itemize}
        \item Stops service
        \item Starts again
        \item Drops active connections
    \end{itemize}
\end{itemize}

Example:

\begin{highlightbox}
\texttt{\$ sudo systemctl reload apache2}
\end{highlightbox}

\subsection*{\faPowerOff\ Enable at Boot}

\begin{highlightbox}
\texttt{\$ sudo systemctl enable apache2}
\end{highlightbox}

This creates a symlink in:

\begin{highlightbox}
\texttt{/etc/systemd/system/}
\end{highlightbox}

\subsection*{\faBan\ Disable at Boot}

\begin{highlightbox}
\texttt{\$ sudo systemctl disable apache2}
\end{highlightbox}

\subsection*{\faSkull\ Mask a Service}

Masking = completely block it from starting.

\begin{highlightbox}
\texttt{\$ sudo systemctl mask apache2}
\end{highlightbox}

Unmask:

\begin{highlightbox}
\texttt{\$ sudo systemctl unmask apache2}
\end{highlightbox}

\begin{infobox}
\arrowicon\ More explanation: \href{https://www.geeksforgeeks.org/linux-unix/how-to-mask-a-systemd-unit-in-linux/}{\texttt{geeksforgeeks-mask-a-systemd-unit-in-linux}}
\end{infobox}

\section*{5 \faFolderOpen\ Where Are Unit Files Stored?}

Systemd has priority layers.

\subsection*{\faDatabase\ Main Directories}

\begin{infobox}
\begin{tabular}{ll}
\textbf{Directory} & \textbf{Purpose} \\
\texttt{/usr/lib/systemd/system/} & Default package units \\
\texttt{/lib/systemd/system/} & (Debian equivalent) \\
\texttt{/etc/systemd/system/} & Custom \& overrides \\
\end{tabular}
\end{infobox}

\subsection*{Priority Order}

Highest priority:

\begin{highlightbox}
\texttt{/etc/systemd/system/}
\end{highlightbox}

So if you edit there — it overrides the default.

\section*{6 \faFileCode\ Anatomy of a Service File}

Let's inspect:

\begin{highlightbox}
\texttt{\$ systemctl cat apache2.service}
\end{highlightbox}

Typical structure:

\begin{highlightbox}
\begin{verbatim}
[Unit]
Description=...
After=network.target

[Service]
Type=forking
ExecStart=/usr/sbin/apachectl start
ExecReload=/usr/sbin/apachectl graceful
ExecStop=/usr/sbin/apachectl stop

[Install]
WantedBy=multi-user.target
\end{verbatim}
\end{highlightbox}

\subsection*{\faList\ [Unit] Section}

Defines metadata and dependencies.

Common options:

\begin{itemize}
    \item Description
    \item After=
    \item Requires=
\end{itemize}

\subsection*{\faCog\ [Service] Section}

How the service runs.

Important fields:

\begin{infobox}
\begin{tabular}{ll}
\textbf{Option} & \textbf{Meaning} \\
ExecStart & Command to start \\
ExecStop & Command to stop \\
Restart & Restart policy \\
User & Run as user \\
Type & simple, forking, oneshot \\
\end{tabular}
\end{infobox}

\subsection*{\faPlug\ [Install] Section}

Defines how it integrates into boot.

\begin{highlightbox}
\texttt{WantedBy=multi-user.target}
\end{highlightbox}

Means: Start in normal multi-user mode.

\section*{7 \faCode\ Creating Your Own Service}

Let's create one.

\subsection*{\faBullseye\ Goal: Run a simple script}

Create script:

\begin{highlightbox}
\texttt{\$ sudo nano /usr/local/bin/hello.sh}
\end{highlightbox}

Add:

\begin{highlightbox}
\begin{verbatim}
#!/bin/bash
echo "Hello from systemd!" >> /tmp/hello.log
\end{verbatim}
\end{highlightbox}

Make executable:

\begin{highlightbox}
\texttt{\$ sudo chmod +x /usr/local/bin/hello.sh}
\end{highlightbox}

\subsection*{Create Service File}

\begin{highlightbox}
\texttt{\$ sudo nano /etc/systemd/system/hello.service}
\end{highlightbox}

Add:

\begin{highlightbox}
\begin{verbatim}
[Unit]
Description=My Hello Service

[Service]
Type=oneshot
ExecStart=/usr/local/bin/hello.sh

[Install]
WantedBy=multi-user.target
\end{verbatim}
\end{highlightbox}

\subsection*{Reload systemd:}

\begin{highlightbox}
\texttt{\$ sudo systemctl daemon-reload}
\end{highlightbox}

\begin{warningbox}
\warningicon\ This is IMPORTANT anytime you:
\begin{itemize}
    \item Create
    \item Edit
    \item Delete
\end{itemize}
unit files.
\end{warningbox}

Start service:

\begin{highlightbox}
\texttt{\$ sudo systemctl start hello.service}
\end{highlightbox}

Check result:

\begin{highlightbox}
\texttt{\$ cat /tmp/hello.log}
\end{highlightbox}

\begin{infobox}
Boom \faBolt\\
You just created a service.
\end{infobox}

\section*{8 \faEdit\ Editing \& Overriding Units (The Smart Way)}

Never edit files in \texttt{/lib/systemd/system/}.

Instead:

\begin{highlightbox}
\texttt{\$ sudo systemctl edit apache2}
\end{highlightbox}

This creates:

\begin{highlightbox}
\texttt{/etc/systemd/system/apache2.service.d/override.conf}
\end{highlightbox}

Add:

\begin{highlightbox}
\begin{verbatim}
[Service]
Restart=always
\end{verbatim}
\end{highlightbox}

Save.

Reload daemon:

\begin{highlightbox}
\texttt{\$ sudo systemctl daemon-reload}
\end{highlightbox}

Restart service.

\section*{9 \timericon\ systemd Timers (Modern Scheduling)}

Now we level up.

Timers = replacement for cron.

\subsection*{Why Use systemd Timers Instead of Cron?}

\begin{itemize}
    \item Integrated with services
    \item Better logging
    \item More control
    \item Dependency aware
    \item Can use calendar expressions
\end{itemize}

\begin{infobox}
\faInfoCircle\ Note that systemd is more complex than Cron and requires systemd to work on your system.
\end{infobox}

\subsection*{\faClock\ Timer Structure}

A timer requires:

\begin{enumerate}
    \item A \texttt{.service} file
    \item A \texttt{.timer} file
\end{enumerate}

\section*{10 \faEnvelope\ Example: Send Message to All Users}

We'll simulate the example from the referenced video.

\subsection*{Create Script}

\begin{highlightbox}
\texttt{\$ sudo nano /usr/local/bin/wall-msg.sh}
\end{highlightbox}

Add:

\begin{highlightbox}
\begin{verbatim}
#!/bin/bash
wall "System reminder: Stay awesome!"
\end{verbatim}
\end{highlightbox}

Make executable:

\begin{highlightbox}
\texttt{\$ sudo chmod +x /usr/local/bin/wall-msg.sh}
\end{highlightbox}

\subsection*{Create Service File}

\begin{highlightbox}
\texttt{\$ sudo nano /etc/systemd/system/wall-msg.service}
\end{highlightbox}

\begin{highlightbox}
\begin{verbatim}
[Unit]
Description=Wall Message Service

[Service]
Type=oneshot
ExecStart=/usr/local/bin/wall-msg.sh
\end{verbatim}
\end{highlightbox}

\subsection*{Create Timer File}

\begin{highlightbox}
\texttt{\$ sudo nano /etc/systemd/system/wall-msg.timer}
\end{highlightbox}

\begin{highlightbox}
\begin{verbatim}
[Unit]
Description=Run wall message every minute

[Timer]
OnCalendar=*-*-* *:*:00
Persistent=true

[Install]
WantedBy=timers.target
\end{verbatim}
\end{highlightbox}

\subsection*{Reload daemon:}

\begin{highlightbox}
\texttt{\$ sudo systemctl daemon-reload}
\end{highlightbox}

Enable timer:

\begin{highlightbox}
\texttt{\$ sudo systemctl enable wall-msg.timer}
\end{highlightbox}

Start timer:

\begin{highlightbox}
\texttt{\$ sudo systemctl start wall-msg.timer}
\end{highlightbox}

Check timers:

\begin{highlightbox}
\texttt{\$ systemctl list-timers}
\end{highlightbox}

\begin{infobox}
�� You are now scheduling tasks.
\end{infobox}

\subsection*{\faCalendar\ Understanding OnCalendar}

Examples:

\begin{infobox}
\begin{tabular}{ll}
\textbf{Expression} & \textbf{Meaning} \\
\texttt{daily} & Once per day \\
\texttt{weekly} & Weekly \\
\texttt{Mon *-*-* 09:00:00} & Mondays 9AM \\
\texttt{*-*-01 00:00:00} & First day of month \\
\end{tabular}
\end{infobox}

Test calendar expression:

\begin{highlightbox}
\texttt{\$ systemd-analyze calendar "Mon *-*-* 09:00:00"}
\end{highlightbox}

\subsection*{\faSlidersH\ Other Timer Options}

\begin{infobox}
\begin{tabular}{ll}
\textbf{Option} & \textbf{Meaning} \\
OnBootSec= & After boot \\
OnUnitActiveSec= & After last run \\
Persistent=true & Run missed jobs after reboot \\
\end{tabular}
\end{infobox}

\section*{11 \faBalanceScale\ Timer vs Cron Comparison}

\begin{infobox}
\begin{tabular}{lll}
\textbf{Feature} & \textbf{Cron} & \textbf{systemd Timer} \\
Logging & Basic & journalctl \\
Dependencies & No & Yes \\
Missed Runs & No & Yes (Persistent) \\
Integration & Separate & Native \\
\end{tabular}
\end{infobox}

\section*{12 \faBug\ Debugging Services \& Timers}

Logs:

\begin{highlightbox}
\texttt{\$ journalctl -u wall-msg.service}
\end{highlightbox}

Live logs:

\begin{highlightbox}
\texttt{\$ journalctl -f}
\end{highlightbox}

Check failures:

\begin{highlightbox}
\texttt{\$ systemctl --failed}
\end{highlightbox}

\section*{13 \faGlobe\ Real World Use Cases}

\begin{itemize}
    \item Backup scripts
    \item Log cleanup
    \item Health checks
    \item Auto-updates
    \item Dev server restarts
    \item IoT tasks
\end{itemize}

\section*{14 \warningicon\ Common Mistakes}

\begin{warningbox}
- Forgetting \texttt{daemon-reload}\\
- Editing \texttt{/lib/systemd/system/} directly\\
- Wrong file permissions\\
- Forgetting \texttt{WantedBy=}
\end{warningbox}

\vspace{1em}
\begin{center}
\color{teamred}\rule{0.8\textwidth}{1pt}
\end{center}

\section*{\faBrain\ Final Mental Model}

\begin{itemize}
    \item Service = What to run
    \item Timer = When to run
    \item Unit file = Configuration blueprint
    \item systemctl = Control center
\end{itemize}

\section*{\rocketicon\ You Now Know}

\begin{itemize}
    \item How systemd works
    \item How to control services
    \item How to create units
    \item How to override safely
    \item How to schedule with timers
    \item How to debug issues
\end{itemize}

\begin{infobox}
\begin{center}
You are no longer a Linux user.\\
\large\textbf{You're becoming a Linux operator.} \rocketicon
\end{center}
\end{infobox}

% Updated Footer
\vspace{\fill}
\begin{center}
\color{muted}
All rights reserved for \texttt{TogetherTeam} [NOT For Commercial Use] !!\\
contact the author via email \href{mailto:khaledmhfz2004@gmail.com}{\texttt{khaledmhfz2004@gmail.com}}
\end{center}

\end{document}